\section{远征的道路比版图更长}
\label{sec:sui-frontiers}

开皇十八年对高句丽的水陆出兵，在疾疫、风浪与补给困难中撤回，已经说明辽东不是可以用一次诏令纳入朝廷秩序的边缘。大业八年的大规模再征，辽东城久攻不下，渡辽水的部队覆败；两年后的第三次出兵虽因高句丽送还斛斯政而撤军，也没有消除双方竞争。三次行动的兵数、伤亡和战场细节主要保存于隋方战史，不能当作可核算的实数；但连续调兵、补给受阻与撤回的序列足以显示，帝国动员的边界正落在道路、气候、疾病和对手的城防之上。%
\footnote{三次征高句丽见 ST-598-GOGURYEO-FIRST-CAMPAIGN、ST-612-GOGURYEO-INVASION、ST-614-GOGURYEO-THIRD，依据《隋书·高祖纪下》《隋书·炀帝纪下》《隋书·东夷传·高丽》及《资治通鉴》隋纪。高句丽方面的连续材料缺环，伤亡和船数不宜照录。}

609年炀帝至河西，隋军击吐谷浑，并置西海、河源等郡。这里也应避免把置郡理解为高原已被均质地纳入州县：军队可以推进，行政名号可以写进地理志，驻防、补给、道路与当地政治却未必随之稳定。河西—青海通路的争夺表明，隋的边疆政策既关乎威望，也关乎能否控制连接关中、西域和高原的资源节点。%
\footnote{吐谷浑战事及置郡见 ST-609-TUYUHUN-CAMPAIGN，依据《隋书·炀帝纪上》《隋书·吐谷浑传》《隋书·地理志》。郡治的持续时间和实际控制范围须由遗址、文书另作检验。}

因此，隋的边疆并不构成一个从中心向外自然展开的圆圈。辽东的水陆运粮、河西的高原通道和北方突厥的压力，都迫使朝廷在同一批兵员、粮食与官吏之间作选择。后来隋末的危机并非由一场远征单独造成；不过，这些远征确实把征发的成本集中地显露出来，并使地方对中央能否兑现命令与供给的信任更易动摇。
